\documentclass[11pt,a4paper]{article}
\usepackage[utf8]{inputenc}
\usepackage[T1]{fontenc}
\usepackage{amsfonts}
\usepackage[french]{babel}
\frenchbsetup{StandardLists=true} % à inclure si on utilise \usepackage[francais]{babel}
\usepackage{enumitem}
\usepackage{amssymb}
\usepackage{amsmath}
\usepackage[left=2cm,right=2cm,top=2cm,bottom=2cm]{geometry}
\usepackage{multirow}
\usepackage[dvipsnames]{xcolor}

\usepackage[T1]{fontenc}
\usepackage{graphicx}

\usepackage{setspace}
\singlespacing

\usepackage{pstricks}
\usepackage{xkeyval}
\usepackage{pst-labo}


\usepackage{multicol}
\usepackage{caption}
\usepackage{fancyhdr}
\pagestyle{fancy}
\renewcommand\headrulewidth{1pt}
\fancyhead[L]{Lycée Denis Diderot }
\fancyhead[R]{Classe de 1 Bac Pro }
\setcounter{page}{3}\fancyfoot[C]{page \thepage}
\fancyfoot[R]{Année 2024-2025}
\fancyfoot[L]{Alexis RUIZ}
\usepackage{multirow,tabularx,booktabs}
\usepackage[tikz]{bclogo} 
\usepackage{circuitikz}

\newcommand{\cadreblanc}[1]{%
\medskip\framebox[\linewidth][c]{%
\rule{0mm}{#1mm}}\par
\medskip}

\usepackage{awesomebox}
\setlength{\aweboxvskip}{0mm}
\usepackage{pifont}
\usepackage{chemmacros}
\chemsetup{modules=all}
\setlength{\columnseprule}{.25mm}

\renewcommand{\arraystretch}{2}
\newcommand*\acocher{\quitvmode{\fboxsep0pt \fboxrule0.8pt \fbox{\vrule height1.8ex depth.1ex width0pt \vrule height0pt depth0pt width1.9ex }}}

\newcolumntype{C}[1]{>{\centering\arraybackslash }b{#1}}

\newcommand{\code}{\awesomebox[black]{2pt}{

\includegraphics[scale=0.15]{Chapitre 1 - Energie et puissance électrique/images/frame.png}
}{black}}

\newcommand{\moteur}{\awesomebox[black]{2pt}{

\includegraphics[scale=1]{images/moteur.jpg} 
}{black}}

\newcommand{\defbox}{\awesomebox[black]{2pt}{
\includegraphics[scale=0.5]{Chapitre 1 - Energie et puissance électrique/images/appel exam.png}}{black}}

\newcommand{\defconclusion}{\awesomebox[black]{2pt}{\rotatebox{90}{\Large{\textbf{Conclusion}}}}{black}}

  \newenvironment{itemize*}%  % listes compactes
    {\begin{itemize}%
      \setlength{\itemsep}{0pt}%
      \setlength{\parskip}{0pt}}%
    {\end{itemize}}

\frenchbsetup{StandardLists=true}



\begin{document}


\flushleft \textbf{NOM : \hfill Classe : \hspace{3cm} ~~\\[2mm]
Prénom:\\[0,4cm]
}
\hrule\
\\[0,2cm]

\begin{center}
 \LARGE{TP1 - Comment évolue la pression en fonction du volume ? }
\end{center}

\vspace{3 mm}

\normalsize

\hrule\
\\[0,2cm]
   \begin{center}
   \textbf{Le soin et la rédaction interviendront dans la notation.}
   \end{center} 

\vspace{4mm}



\fcolorbox{black}{white}{\textbf{1. Expérience}}  \\
\vspace{-15mm}
 \begin{enumerate}[label*=1.\arabic*.]
\item 
\noindent%
\begin{minipage}{.4\textwidth}%
 \begin{bclogo}[couleur=white ,  arrondi =0.1 , logo = \bcinfo, epBarre =2]{ Matériel }

\begin{itemize*}
    \item 1 seringue de 100 mL graduée tous les 10 mL.
    \item 1 tuyaux souple de 3mL
    \item 1 capteur de pression
    \item 1 ordinateur avec Atelier Scientifique Jeulin
    \item 1 interface Primo
\end{itemize*}
\end{bclogo}

\end{minipage}%
\hfill
\begin{minipage}{.5\textwidth}%

\begin{center}
\includegraphics[width=\textwidth]{Chapitre 3 - Pression et forces préssantes/images/Picto_primo.png}
 
\end{center}
\end{minipage}%



Il suffit de raccorder l'extrémité de la seringue au tuyau souple du capteur de pression.


On veut mesurer la pression en fonction du volume occupé par le gaz dans la seringue. \\ Pour cela, on fera une série de 5 mesures de 103 à 53 mL, et la pression correspondante sera enregistrée par le logiciel. On doit tenir compte du volume du tuyau souple.

\vfill 
    
\defbox{\fbox{
\begin{minipage}{8cm}{

Appeler M Ruiz pour qu'il vérifie le montage.}

\end{minipage}
}}

\vfill

\item Paramétrage de l'acquisition : 
\noindent%
\begin{minipage}{.85\textwidth}%

\includegraphics{Chapitre 3 - Pression et forces préssantes/images/image1.png} 

\end{minipage}%

\newpage 

\item Clique sur l'icône \raisebox{-1.2
ex} {\includegraphics[scale=0.5]{Chapitre 3 - Pression et forces préssantes/images/icone.png}} pour afficher la fenêtre de paramétrage de l'acquisition. \\ On veut visualiser $p = f(V)$.\\[3mm]

\noindent%
\begin{minipage}{.45\textwidth}%

\'A l'aide de la souris faire glisser l'adaptateur pression sur l'axe des ordonnés. \\[2mm]

\includegraphics[scale=0.85]{Chapitre 3 - Pression et forces préssantes/images/image2.png} 

\end{minipage}%
\hfill
\begin{minipage}{.45\textwidth}%

\'A l'aide de la souris faire glisser le clavier sur l'axe des abscisses. \\[2mm]
\includegraphics[scale=0.85]{Chapitre 3 - Pression et forces préssantes/images/image3.png} 

\end{minipage}%

\vfill 


\noindent%
\begin{minipage}{.45\textwidth}%

\textbf{Capteur de pression :}  \\[2mm]

\includegraphics[scale=0.95]{Chapitre 3 - Pression et forces préssantes/images/image4.png} 

\end{minipage}%
\hfill
\begin{minipage}{.45\textwidth}%

\textbf{Clavier : } \\[2mm]

Cliquer sur l'onglet d'option Grandeur et compléter : \\[2mm]

\includegraphics[scale=0.95]{Chapitre 3 - Pression et forces préssantes/images/image5.png} 

\end{minipage}%


\vfill 


\defbox{\fbox{
\begin{minipage}{9cm}{

Appeler M Ruiz pour qu'il vérifie le paramétrage.}

\end{minipage}
}}

\vfill 

\item Acquisition : 

\begin{itemize}
    \item 


\noindent%
\begin{minipage}{.85\textwidth}%

\includegraphics{Chapitre 3 - Pression et forces préssantes/images/image6.png} 

\end{minipage}%

\item Régler le volume de la seringue  sur 100 mL, renseigner la boite de dialogue  et cliquer sur le bouton.  \raisebox{-2
ex} {\includegraphics[]{Chapitre 3 - Pression et forces préssantes/images/icone2.jpg}}






\newpage

\item Pousser le piston de la seringue de 10 mL , renseigner le nouveau volume de la seringue puis cliquer sur OK Suivant.  

\item Faire de même jusqu'à 53 mL.

\item Cliquer sur \raisebox{-2
ex}{
\includegraphics[]{Chapitre 3 - Pression et forces préssantes/images/icone3.jpg} }pour arrêter l'acquisition. 

\item Enregistrer le fichier (Enregistrer sous).

\item La courbe peut être agrandie en cliquant sur l’icône Échelle automatique \raisebox{-1.5
ex}{\includegraphics[]{Chapitre 3 - Pression et forces préssantes/images/icone4.jpg} }.

\end{itemize}
\end{enumerate}




\fcolorbox{black}{white}{\textbf{2. Modélisation}}  \\

\begin{itemize}
    \item Dans l'onglet Tableau créer une nouvelle grandeur p$\times$V unité SI

    \item Calculer le produit p$\times$V pour les 5 valeurs. 

    \item Commenter les résultats du produit p$\times$V. \\[8mm]

    \dotfill

    \vspace{5mm}
    On veut modéliser $p=f(V)$ : \\[3mm]
    
  \item Clique sur l'icône   \raisebox{-1.5
ex}{\includegraphics[]{Chapitre 3 - Pression et forces préssantes/images/icone5.jpg} },       et demander la modélisation pour la courbe obtenue. 


  \item Cliquer sur l'onglet modélisation automatique et saisir pm = constante/V dans le champ de texte correspondant.
  
  \item Cliquer ensuite sur Modéliser et relever la valeur de la constante. \\[8mm]
  
  \dotfill 

\vfill 


\defbox{\fbox{
\begin{minipage}{9cm}{

Appeler M Ruiz pour qu'il vérifie vos résultats.}

\end{minipage}
}}

\vfill 
    
\end{itemize}

\vfill 

\fcolorbox{black}{white}{\textbf{3. Exploitation}}  \\[3mm]

\begin{itemize}
    \item 


  Quel type de courbe obtient-on ? \\[3mm]

 \dotfill 
 
  Comment évoluent les grandeurs pression et volume au cours de cette expérience ? \\[3mm]

 \dotfill 

 \vspace{3mm}

 \dotfill 

 
 
\item Parmi ces affirmations, cocher celles qui sont justes : \\[3mm]

 \begin{itemize}[label=$\square$]
 

  \item A température constante, quand la pression d'un gaz augmente, son volume augmente. 
  
  \item A température constante, le volume d'un gaz et sa pression sont des grandeurs       proportionnelles. 
  
   \item A température constante, le volume d’un gaz est inversement proportionnel à la pression        qu’il reçoit. 
   
   \item A température constante, la pression d'un gaz évolue en fonction du carré de son volume. 

   \end{itemize}
 \end{itemize}
\end{document}